%%%%%%%%%%%%%%%%%%%%%%%%%%%%%%%%%
%
%       EXERCÍCIO
%
%%%%%%%%%%%%%%%%%%%%%%%%%%%%%%%%%

\ifdefstring{\atividade}{prova}{%
    \ifdefstring{\modo}{objetivo}{%
        \renewcommand{\valorquestao}{\ValorQObj\ ponto}
    }{%
        \renewcommand{\valorquestao}{\ValorQDisc\ pontos}
    }%
}%

\begin{exercicioBanco}[\valorquestao]
Sabendo que o gráfico a seguir representa a função \(f(x) = |x-2| + |x+3|\), determine o valor de \(a+b+c\).
\begin{center}
    \begin{tikzpicture}[scale=0.5]
        % Eixos
        \draw[->] (-4,0) -- (3,0) node[right] {\(x\)};
        \draw[->] (0,-.5) -- (0,7) node[above] {\(y\)};

        % Gráfico da função f(x) = |x-2| + |x+3|
        \draw[domain=-4:-3, thick, blue] plot (\x, {-2*\x -1}); % x < -1
        \draw[domain=-3:2, thick, blue] plot (\x, {5});    % x ≥ -1
        \draw[domain=2:3, thick, blue] plot (\x, {2*\x + 1});    % x ≥ -1

        % Marcação do vértice (x = -1)
        \draw[dashed] (-3,0) -- (-3,5);
        \node[below] at (-3,0) {\small\(a\)};
        \draw[dashed] (2,0) -- (2,5);
        \node[below] at (2,0) {\small\(b\)};
        \node[above right] at (0,5) {\small\(c\)};
    \end{tikzpicture}
\end{center}


% Define as alternativas
\newcommand{\alternativas}{%
    \begin{center}
        \begin{tabularx}{\textwidth}{XXXXX}
            (a) \(-7\). &
            (b) \(-6\). &
            (c) \(4\). &
            (d) \(6\). &
            (e) \(10\).
        \end{tabularx}
    \end{center}
}

% Define a resposta correta
\newcommand{\resposta}{C}

% Lógica condicional para exibição
\ifdefstring{\atividade}{lista}{%
    \alternativas
    \vspace{0.5em}
    
    \noindent\textbf{Resposta:} letra \textbf{\resposta}.
}{%
    \ifdefstring{\modo}{objetivo}{%
        \alternativas
    }{%
        % modo = discursiva → não mostra alternativas
    }
}
\end{exercicioBanco}

